% ============================================================
%  Corriges Chapitre 4 : Systemes lineaires et programmation lineaire
% ============================================================
\clearpage
\chapter{Systèmes linéaires et programmation linéaire}
\chaptermark{Systèmes linéaires et programmation linéaire}

\begin{corrige}{4.1}
\textbf{Par substitution.} La deuxième équation donne $x = y + 1$ ; en reportant dans la première : $3(y + 1) + 2y = 13$, soit $5y + 3 = 13$, donc $5y = 10$, $y = 2$ et $x = 3$.
\par\smallskip
\textbf{Par combinaison.} Calculons $(\text{L}_{1}) - 3(\text{L}_{2})$ : les termes en $x$ s'éliminent ($3x - 3x = 0$) et il reste $2y + 3y = 13 - 3$, soit $5y = 10$ et $y = 2$, puis $x = y + 1 = 3$.
\par\smallskip
Dans les deux cas $S = \{(3;\ 2)\}$ ; la vérification donne $3 \times 3 + 2 \times 2 = 13$ et $3 - 2 = 1$. Les deux méthodes conduisent toujours au même couple : le choix n'est qu'une affaire de calcul.
\end{corrige}

\begin{corrige}{4.2}
\begin{enumerate}
\item Les déterminants valent $D = 2 \times (-1) - 3 \times 4 = -14$, $D_{x} = 12 \times (-1) - 3 \times 10 = -42$ et $D_{y} = 2 \times 10 - 12 \times 4 = -28$. Comme $D \neq 0$ : $x = \dfrac{D_{x}}{D} = \dfrac{-42}{-14} = 3$ et $y = \dfrac{D_{y}}{D} = \dfrac{-28}{-14} = 2$. La solution est $S = \{(3;\ 2)\}$ ; vérification : $6 + 9 = 12$ et $12 - 2 = 10$.
\item $(S_{1})$ : le déterminant du système vaut $D = 2 \times 3 - 1 \times 6 = 0$, avec $D_{x} = 5 \times 3 - 1 \times 15 = 0$ et $D_{y} = 2 \times 15 - 5 \times 6 = 0$ : les trois déterminants sont nuls, les deux équations sont équivalentes (la deuxième est le triple de la première) et le système a une infinité de solutions, $S = \{(x;\ 5 - 2x),\ x \in \R\}$. $(S_{2})$ : $D = 0$ encore, mais $D_{x} = 5 \times 3 - 1 \times 8 = 7 \neq 0$ : le système n'a aucune solution, $S = \varnothing$ (les droites sont parallèles et distinctes).
\end{enumerate}
\end{corrige}

\begin{corrige}{4.3}
On élimine $x$ : $(\text{L}_{2}) - 2(\text{L}_{1})$ donne $-3y - z = -9$, et $(\text{L}_{3}) - (\text{L}_{1})$ donne $y - 2z = -4$. On élimine $y$ : calculons $3(\text{L}_{3}') + (\text{L}_{2}')$, où $\text{L}_{3}'$ désigne la nouvelle troisième ligne : $3y - 6z = -12$ puis $3y - 6z - 3y - z = -12 - 9$, soit $-7z = -21$ et $z = 3$. Remontée : $y - 2 \times 3 = -4$ donne $y = 2$, puis $x + 2 + 3 = 6$ donne $x = 1$.
\par\smallskip
La solution est $S = \{(1;\ 2;\ 3)\}$. Contrôle dans les trois équations initiales : $1 + 2 + 3 = 6$, $2 - 2 + 3 = 3$ et $1 + 4 - 3 = 2$ : le triplet vérifie les trois équations.
\end{corrige}

\begin{corrige}{4.4}
La droite d'équation $x + y = 6$ passe par $(0;\ 6)$ et $(6;\ 0)$ ; la droite d'équation $2x - y = 4$ passe par $(0;\ -4)$ et $(2;\ 0)$. Le test de l'origine : $0 + 0 = 0 \leqslant 6$ est vérifié et $2 \times 0 - 0 = 0 \leqslant 4$ aussi : pour chaque inéquation, on garde le demi-plan contenant $O$ et on hachure l'autre ; les axes $x \geqslant 0$ et $y \geqslant 0$ limitent le dessin au premier quadrant.
\par\smallskip
Le domaine solution est un quadrilatère de sommets $O(0;\ 0)$, $(2;\ 0)$, l'intersection des deux droites qui vérifie $x + 2x - 4 = 6$, soit $x = \dfrac{10}{3}$ et $y = \dfrac{8}{3}$ (point $\left(\dfrac{10}{3};\ \dfrac{8}{3}\right)$), et $(0;\ 6)$ : chaque sommet est le point où deux frontières se croisent, bord inclus puisque les inégalités sont larges.
\end{corrige}

\begin{corrige}{4.5}
\begin{enumerate}
\item Les quantités sont positives et les disponibilités se traduisent par :
\[
\begin{cases}
x \geqslant 0,\quad y \geqslant 0\\
2x + 3y \leqslant 60\\
x + 2y \leqslant 34
\end{cases}
\]
\item Le bénéfice hebdomadaire vaut $Z = 2\,500x + 4\,000y$ (en francs).
\item Le domaine des contraintes a pour sommets $O(0;\ 0)$, $A(30;\ 0)$ (sur la droite de planches avec $y = 0$), $E(18;\ 8)$ et $B(0;\ 17)$ (sur la droite d'heures avec $x = 0$). Le point $E$ est l'intersection des deux droites de contrainte : de $x = 34 - 2y$ reporté dans $2x + 3y = 60$, on tire $68 - 4y + 3y = 60$, donc $y = 8$ et $x = 18$.
\item On calcule $Z$ à chaque sommet :
\[
Z(O) = 0, \qquad Z(A) = 2\,500 \times 30 = 75\,000, \qquad Z(B) = 4\,000 \times 17 = 68\,000,
\]
\[
Z(E) = 2\,500 \times 18 + 4\,000 \times 8 = 45\,000 + 32\,000 = 77\,000 .
\]
Le maximum vaut $77\,000$ F au sommet $E$ : la menuiserie doit fabriquer $18$ chaises et $8$ tables par semaine. Contrôle des ressources : $2 \times 18 + 3 \times 8 = 60$ planches et $18 + 2 \times 8 = 34$ heures : les deux ressources sont utilisées en totalité, comme il se doit au sommet qui croise les deux contraintes.
\end{enumerate}
\end{corrige}
